% from ega2-1-fr.tex
\oldpage[II]{88}

puisque $\mathcal{F}(n)$ est un foncteur exact en $\mathcal{F}$~; d'où la
suite exacte de cohomologie
\[
  0\longrightarrow\Gamma(X,\mathcal{J}(n))\longrightarrow
  \Gamma(X,\mathcal{L}^{\otimes n})\longrightarrow
  \Gamma(X,\mathcal{L}'^{\otimes n})\xrightarrow{\partial}
  \mathrm{H}^1(X,\mathcal{J}(n))
\]
qui associe en particulier à $g$ un élément
$\partial g\in\mathrm{H}^1(X,\mathcal{J}(n))$.

Notons que puisque $\mathcal{J}^2=0$, $\mathcal{J}$ peut être considéré
comme un $\mathcal{O}_Z$-Module quasi-cohérent et l'on a, pour tout $k$,
$\mathcal{L}'^{\otimes k}\otimes_{\mathcal{O}_Z}\mathcal{J}(n)
=\mathcal{J}(n+k)$~; pour toute section
$s\in\Gamma(X,\mathcal{L}'^{\otimes k})$ la multiplication tensorielle par
$s$ est donc un homomorphisme $\mathcal{J}(n)\to\mathcal{J}(n+k)$ de
$\mathcal{O}_Z$-Modules, qui donne par suite un homomorphisme
$\mathrm{H}^i(X,\mathcal{J}(n))\to
\mathrm{H}^i(X,\mathcal{J}(n+k))$ de groupes de cohomologie.

Cela étant, nous allons voir que l'on a
\begin{equation}
\label{II.4.5.13.2-fr}
  g^{\otimes m}\otimes\partial g=0
\tag{4.5.13.2}
\end{equation}
pour un $m>0$ assez grand. En effet, $Z_g$ est un ouvert affine de $Z$ et
on a donc $\mathrm{H}^1(Z_g,\mathcal{J}(n))=0$ quand $\mathcal{J}(n)$ est
considéré comme un $\mathcal{O}_Z$-Module
(\hyperref[I.5.1.9.2-fr]{I, 5.1.9.2}). En particulier, si on pose
$g'=g|Z_g$, et si on considère son image par l'application
$\partial:\Gamma(Z_g,\mathcal{L}'^{\otimes n})\to
\mathrm{H}^1(Z_g,\mathcal{J}(n))$, on a $\partial g'=0$. Pour expliciter
cette relation, observons qu'en dimension $1$ la cohomologie d'un faisceau
de groupes abéliens est la même que sa cohomologie de Čech (G, II, 5.9)~;
pour former $\partial g$, il faut donc considérer un recouvrement ouvert
$(U_\alpha)$ assez fin de $X$, que l'on peut supposer fini et formé
d'ouverts affines, prendre pour chaque $\alpha$ une section
$g_\alpha\in\Gamma(U_\alpha,\mathcal{L}^{\otimes n})$ dont l'image
canonique dans $\Gamma(U_\alpha,\mathcal{L}'^{\otimes n})$ est
$g|U_\alpha$, et considérer la classe du cocycle
$(g_{\alpha|\beta}-g_{\beta|\alpha})$, $g_{\alpha|\beta}$ étant la
restriction de $g_\alpha$ à $U_\alpha\cap U_\beta$ (ce cocycle étant à
valeurs dans $\mathcal{J}(n)$). On peut en outre supposer que $\partial g'$
est calculé de la même manière à l'aide du recouvrement formé des
$U_\alpha\cap Z_g$ et des restrictions $g_\alpha|(U_\alpha\cap Z_g)$ (en
remplaçant au besoin $(U_\alpha)$ par un recouvrement plus fin)~; la
relation $\partial g'=0$ signifie alors qu'il existe pour chaque $\alpha$
une section $h_\alpha\in\Gamma(U_\alpha\cap Z_g,\mathcal{J}(n))$ telle que
\[
  (g_{\alpha|\beta}-g_{\beta|\alpha})
  |(U_\alpha\cap U_\beta\cap Z_g)
  =h_{\alpha|\beta}-h_{\beta|\alpha}
\]
en désignant par $h_{\alpha|\beta}$ la restriction de $h_\alpha$ à
$U_\alpha\cap U_\beta\cap Z_g$ (G, II, 5.11). Cela étant, il y a un entier
$m>0$ tel que $g^{\otimes m}\otimes h_\alpha$ soit la restriction à
$U_\alpha\cap Z_g$ d'une section
$t_\alpha\in\Gamma(U_\alpha,\mathcal{J}(n+nm))$ pour tout $\alpha$
(\hyperref[I.9.3.1-fr]{I, 9.3.1})~; on a donc
\[
  g^{\otimes m}\otimes(g_{\alpha|\beta}-g_{\beta|\alpha})
  =t_{\alpha|\beta}-t_{\beta|\alpha}
\]
pour tout couple d'indices, ce qui prouve
(\hyperref[II.4.5.13.2-fr]{4.5.13.2}).

Remarquons d'autre part que si
$s\in\Gamma(X,\mathcal{O}_Z(p))$,
$t\in\Gamma(X,\mathcal{O}_Z(q))$, on a, dans le groupe
$\mathrm{H}^1(X,\mathcal{J}(p+q))$
\begin{equation}
\label{II.4.5.13.3-fr}
  \partial(s\otimes t)=(\partial s)\otimes t+s\otimes(\partial t).
\tag{4.5.13.3}
\end{equation}

En effet, on peut encore, pour calculer les deux membres, considérer un
recouvrement ouvert $(U_\alpha)$ de $X$, pour chaque $\alpha$ une section
$s_\alpha\in\Gamma(U_\alpha,\mathcal{O}_X(p))$ (resp.
$t_\alpha\in\Gamma(U_\alpha,\mathcal{O}_X(q))$ dont l'image canonique dans
$\Gamma(U_\alpha,\mathcal{O}_Z(p))$ (resp.
$\Gamma(U_\alpha,\mathcal{O}_Z(q))$ soit $s|U_\alpha$ (resp.
$t|U_\alpha$)~; la relation
(\hyperref[II.4.5.13.3-fr]{4.5.13.3}) résulte alors des relations
\[
  (s_{\alpha|\beta}\otimes t_{\alpha|\beta})
  -(s_{\beta|\alpha}\otimes t_{\beta|\alpha})
  =(s_{\alpha|\beta}-s_{\beta|\alpha})\otimes t_{\alpha|\beta}
  +s_{\beta|\alpha}\otimes(t_{\alpha|\beta}-t_{\beta|\alpha})
\]
avec les mêmes notations que ci-dessus. Par récurrence sur $k$, on a donc
\begin{equation}
\label{II.4.5.13.4-fr}
  \partial(g^{\otimes k})=(kg^{\otimes(k-1)})\otimes(\partial g)
\tag{4.5.13.4}
\end{equation}

